\begin{register}{H}{INT\_ENABLE\_REG}{0x0104}\label{regdesc:INTENABLEREG}
\regfield{INT\_ENABLE[31:8]}{24}{8}{{0x00000000}}%
\regfield{(reserved)}{1}{7}{{0}}%
\regfield{INT\_ENABLE[6:5]}{2}{5}{{0x0}}%
\regfield{(reserved)}{2}{3}{{0x0}}%
\regfield{INT\_ENABLE[2:1]}{2}{1}{{0x0}}%
\regfield{(reserved)}{1}{0}{{0}}%
\reglabel{Reset}\regnewline%
\begin{regdesc}
\begin{reglist}
\item [INT\_ENABLE\lbrack\regindex{n}\rbrack] 设置 \regindex{n} 位以使能 \regindex{n} 外部中断对 CPU 的断言。(R/W)
\begin{itemize}
    \setlength\itemsep{-1em}
    \item 0: 禁用;
    \item 1: 启用;
\end{itemize}
\end{reglist}
\end{regdesc}
\end{register}

\begin{register}{H}{INT\_TYPE\_REG}{0x0108}\label{regdesc:INTTYPEREG}
\regfield{INT\_TYPE[31:8]}{24}{8}{{0x00000000}}%
\regfield{(reserved)}{1}{7}{{0}}%
\regfield{INT\_TYPE[6:5]}{2}{5}{{0x0}}%
\regfield{(reserved)}{2}{3}{{0x0}}%
\regfield{INT\_TYPE[2:1]}{2}{1}{{0x0}}%
\regfield{(reserved)}{1}{0}{{0}}%
\reglabel{Reset}\regnewline%
\begin{regdesc}
\begin{reglist}
\item [INT\_TYPE\lbrack\regindex{n}\rbrack] 设置 \regindex{n} 位以捕获 \regindex{n} 外部中断的上升沿。(R/W)
\begin{itemize}
    \setlength\itemsep{-1em}
    \item 0: 电平类型（信号电平检测）;
    \item 1: 脉冲类型（上升沿检测）;
\end{itemize}
\end{reglist}
\end{regdesc}
\end{register}

\begin{register}{H}{INT\_CLEAR\_REG}{0x010C}\label{regdesc:INTCLEARREG}
\regfield{INT\_CLEAR[31:8]}{24}{8}{{0x00000000}}%
\regfield{(reserved)}{1}{7}{{0}}%
\regfield{INT\_CLEAR[6:5]}{2}{5}{{0x0}}%
\regfield{(reserved)}{2}{3}{{0x0}}%
\regfield{INT\_CLEAR[2:1]}{2}{1}{{0x0}}%
\regfield{(reserved)}{1}{0}{{0}}%
\reglabel{Reset}\regnewline%
\begin{regdesc}
\begin{reglist}
\item [INT\_CLEAR\lbrack\regindex{n}\rbrack] 设置 \regindex{n} 位以清除 \regindex{n} 外部中断的挂起状态。(R/W) \newline 这仅对“脉冲”类型中断有用，因为“电平”类型中断必须在源头清除。 \newline 注意，设置的位必须手动切换回 0。
\begin{itemize}
    \setlength\itemsep{-1em}
    \item 0: 无关紧要;
    \item 1: 清除挂起状态;
\end{itemize}
\end{reglist}
\end{regdesc}
\end{register}

\begin{register}{H}{INT\_EIP\_REG}{0x0110}\label{regdesc:INTEIPREG}
\regfield{INT\_EIP[31:8]}{24}{8}{{0x00000000}}%
\regfield{(reserved)}{1}{7}{{0}}%
\regfield{INT\_EIP[6:5]}{2}{5}{{0x0}}%
\regfield{(reserved)}{2}{3}{{0x0}}%
\regfield{INT\_EIP[2:1]}{2}{1}{{0x0}}%
\regfield{(reserved)}{1}{0}{{0}}%
\reglabel{Reset}\regnewline%
\begin{regdesc}
\begin{reglist}
\item [INT\_EIP\lbrack\regindex{n}\rbrack] 读取 \regindex{n} 位以获取 \regindex{n} 外部中断对 CPU 的挂起状态。(RO) \newline 只有启用且超过阈值的中断才会在此反映。
\begin{itemize}
    \setlength\itemsep{-1em}
    \item 0: 未挂起
    \item 1: 挂起
\end{itemize}
\end{reglist}
\end{regdesc}
\end{register}

\begin{register}{H}{INT\_PRIORITY\_\regindex{n}\_REG (\regindex{n}: 1-2,5-6,8-31)}{0x0114+4*\regindex{n}}\label{regdesc:INTPRIORITYREG}
\regfield{(reserved)}{28}{4}{{0x00000000}}%
\regfield{INT\_PRIORITY\_\regindex{n}}{4}{0}{{0x0}}%
\reglabel{Reset}\regnewline%
\begin{regdesc}
\begin{reglist}
\label{fielddesc:INTPRIORITY}\item [INT\_PRIORITY\_\regindex{n}] 写入一个 4 位值到 \regindex{n} 寄存器以配置 \regindex{n} 外部中断的优先级。(R/W)
\end{reglist}
\end{regdesc}
\end{register}

\begin{register}{H}{INT\_THRESH\_REG}{0x0194}\label{regdesc:INTTHRESHREG}
\regfield{(reserved)}{28}{4}{{0x00000000}}%
\regfield{INT\_THRESH}{4}{0}{{0x0}}%
\reglabel{Reset}\regnewline%
\begin{regdesc}
\begin{reglist}
\label{fielddesc:INTTHRESH}\item [INT\_THRESH] 写入一个 4 位值以配置所有外部中断的全局优先级阈值。(R/W) \newline 所有优先级低于阈值的中断将被屏蔽。 \newline
\end{reglist}
\end{regdesc}
\end{register}
